\subsection{Organización por roles}
\label{sec:backend-rol}

El backend filtra la información por rol (\ac{RBAC}) de forma transversal:
la autenticación y los permisos se gestionan en la app \texttt{usuario}, y
cada app aplica \texttt{get\_queryset} filtrado por el usuario autenticado.
Así, el rol \textbf{estudiante} interactúa con \texttt{cupos},
\texttt{mensajes} y \texttt{evaluaciones}; el rol \textbf{docente} aporta a
\texttt{materias}, \texttt{evaluaciones} (asistencia, justificaciones,
resultados) y \texttt{mensajes}; y la \textbf{secretaría} accede a
\texttt{reportes} (listas de espera, \ac{KPI}) y \texttt{quejas}, con
intervención trazada en hilos escalados. El \textbf{departamento}
configura parámetros globales (preguntas de evaluación, \ac{KPI}) y la app
\texttt{analitica} consolida datos para el motor de soluciones rápidas.

La \figureref{fig:backend} muestra la arquitectura del backend: el servidor
de aplicaciones de \ac{DJ} atiende peticiones \ac{REST} desde el frontend,
persiste en la base de datos y alimenta al motor analítico de forma
asíncrona.
% =====================================================================
%  CAPITULO 1: ARQUITECTURA DE SOFTWARE DEL SGDIC
%  Universidad Nacional de Colombia, Sede Manizales
% =====================================================================

\chapter{Arquitectura de Software del Sistema}
\label{ch:arquitectura}

El \ac{SGDIC} se diseña bajo una arquitectura de tres capas
\textit{(presentación, aplicación y datos)} sobre la plataforma
\ac{DJ}. Esta organización separa claramente la interfaz de usuario, la
lógica de negocio y el acceso a la información, facilitando mantenimiento,
pruebas e incorporación de nuevos requerimientos del modelado \ac{UML}
documentado en el informe de planificación del proyecto.

La figura~\ref{fig:arquitectura-capas} muestra la vista general: cada capa
se comunica con la inferior mediante \ac{API}s de estilo \ac{REST}, y la
capa analítica se nutre de forma asíncrona mediante un motor de extracción
\ac{ETL}.

\begin{figure}[H]
\centering
\adjustbox{max width=\linewidth}{%
\begin{tikzpicture}[node distance=1.5cm and 0.9cm]
  \node[nodoAplicacion, text width=3.7cm, minimum height=1.5cm] (fe) {\textbf{Cap. presentación}\\[2pt]\textit{Frontend}};
  \node[nodoAplicacion, text width=3.7cm, minimum height=1.5cm, right=1.1cm of fe] (be) {\textbf{Cap. aplicación}\\[2pt]\textit{Backend Django}};
  \node[nodoAplicacion, text width=3.7cm, minimum height=1.5cm, right=1.1cm of be] (ce) {\textbf{Cap. datos}\\[2pt]\textit{PostgreSQL}};
  \node[infra, text width=3.1cm, minimum height=1.2cm, below=1.6cm of fe] (etl) {Motor ETL};
  \node[infra, text width=3.1cm, minimum height=1.2cm, right=1.1cm of etl] (dwh) {Repositorio analítico};
  \node[infra, text width=3.1cm, minimum height=1.2cm, right=1.1cm of dwh] (ml) {Motor de soluciones rápidas};
  \draw[flecha] (fe) -- (be);
  \draw[flecha] (be) -- (ce);
  \draw[flechaS] (ce) -- (etl);
  \draw[flechaS] (etl) -- (dwh);
  \draw[flechaS] (dwh) -- (ml);
  \draw[flechaV, dashed] (ml.north) to[bend left=18] (be.south);
\end{tikzpicture}
}%
\caption{Vista de capas de la arquitectura del \ac{SGDIC}}
\label{fig:arquitectura-capas}
\end{figure}

\noindent\begin{figure}[H]
\centering
\adjustbox{max width=\linewidth}{%
\begin{tikzpicture}[
  node distance=1.0cm and 0.9cm,
  rol/.style={compApp, text width=3.1cm, minimum height=1.0cm,
    align=center, inner sep=5pt, font=\small},
  permiso/.style={infra, text width=3.1cm, minimum height=0.9cm,
    align=center, inner sep=5pt, font=\small}
]
  \node[rol, compWeb] (e) {Estudiante};
  \node[rol, right=of e] (d) {Docente};
  \node[rol, right=of d] (s) {Secretaría y\\Departamento};
  \node[permiso, below=1.1cm of e] (nav) {Barra lateral\\responsiva};
  \draw[flecha] (e.south) -- (nav.north);
  \draw[flecha] (d.south) -- (nav.north);
  \draw[flecha] (s.south) -- (nav.north);
\end{tikzpicture}
}%
\caption{Roles y visibilidad de la interfaz del \ac{SGDIC}}
\label{fig:roles-cap}
\end{figure}

\section{Capa de Presentación (Frontend)}
\label{sec:frontend}

La capa de presentación corresponde al \textit{frontend} del \ac{SGDIC},
construido con plantillas \ac{HTML} renderizadas por \ac{DJ} y estilizadas con
\ac{CSS} (Bootstrap) y componentes interactivos en \ac{JS}. La interfaz es
responsiva y sirve tanto a estudiantes como a docentes, secretaría y
departamento, adaptándose al rol autenticado (\ac{RBAC}).

\subsection{Tecnologías y lenguajes}
\label{sec:frontend-tecnologias}

El cliente se construye con pila web estándar renderizada del lado del
servidor (\ac{MVT}): \ac{HTML}5 semántico para la estructura, \ac{CSS}3 con
Bootstrap 5 (\textit{grid}, utilidades y componentes responsivos) para el
estilo y \ac{JS} (ES6+) para la interactividad ligera. El \ac{JS} incluye
jQuery para manipulación DOM y llamadas asíncronas a la \ac{API} \ac{REST},
y Bootstrap JS para componentes dinámicos (modales, tooltips,
off-canvas). No se adopta arquitectura \ac{SPA}; las vistas se renderizan en
el servidor y el navegador recibe marcado ya armado, lo que favorece el
indexado y la carga inicial, manteniendo compatibilidad con entornos con
soporte limitado.

\subsection{Organización por roles}
\label{sec:frontend-rol}

La interfaz se muestra a partir de los tres roles definidos en el
documento de planificación (\cref{fig:roles-cap}):
\begin{itemize}[label=-- , leftmargin=*]
  \item[\textbf{Estudiante}] portal de cupos en tiempo real,
  confirmación/rechazo de horarios, presentación de cuestionarios,
  visualización de faltas y justificaciones, mensajería con docentes y
  autoevaluación histórica.
  \item[\textbf{Docente}] gestión de materias y horarios, registro rápido de
  asistencia (lista o código), aprobación/rechazo de justificaciones,
  envío de mensajes masivos e individuales, consulta de resultados
  consolidados y comparativo histórico.
  \item[\textbf{Secretaría/departamento}] panel de listas de espera,
  intervención en hilos escalados, tableros consolidados de \ac{KPI} y
  exportación de reportes a \ac{CSV} y \ac{PDF}.
\end{itemize}

La adaptación se realiza mediante \ac{RBAC}: la barra lateral y el menú
navegacional sólo muestran los ítems autorizados al rol autenticado, y
cada app del backend filtra su queryset según el usuario. La misma
plantilla se reutiliza para todos los roles; el contenido visible depende
del contexto que el backend inyecta a la vista.

La \figureref{fig:frontend} muestra la arquitectura del cliente: el navegador
consume las vistas renderizadas del servidor y, para interacciones en
tiempo real, delega en notificaciones push y en la API REST del
backend. Los recursos estáticos (\texttt{css}, \texttt{js}, imágenes) se
entregan comprimidos y con caché mediante la configuración de archivos
\texttt{STATICFILES} y \texttt{Whitenoise} (o CDN en despliegue productivo).
\begin{figure}[H]
\centering
\adjustbox{max width=\linewidth}{%
\begin{tikzpicture}[
  node distance=1.2cm and 1.5cm,
  nodo/.style={compApp, text width=3.0cm, minimum height=1.3cm,
    align=center, inner sep=6pt, font=\small}
]
  % Fila 1: cliente y capa de presentacion servida
  \node[nodo, compWeb] (nav) {Navegador del usuario};
  \node[nodo, right=2.6cm of nav] (tmpl) {Plantillas HTML\\(Django)};
  \node[nodo, right=of tmpl] (css) {Bootstrap CSS};
  \node[nodo, below=1.5cm of css] (js) {Bootstrap JS\\y jQuery};
  \node[nodo, compApi, below=1.5cm of tmpl] (api) {API REST del backend};
  \node[nodo, infra, below=1.5cm of js] (static) {Whitenoise / CDN};
  % Flujo principal (azul)
  \draw[flecha] (nav.east) -- node[above, font=\scriptsize]{petición} (tmpl.west);
  \draw[flecha] (tmpl.south) -- node[right, font=\scriptsize]{render} (api.north);
  % Tecnologias del cliente
  \draw[flecha] (tmpl.east) -- (css.west);
  \draw[flecha] (css.south) -- (js.north);
  % Recursos estaticos
  \draw[flechaS, dashed] (static.west) -- ++(-0.8,0) |- (js.south west);
  \draw[flechaS, dashed] (static.north) -- (api.south east);
\end{tikzpicture}
}%
\caption{Arquitectura del frontend del \ac{SGDIC}}
\label{fig:frontend}
\end{figure}

\section{Capa de Aplicación (Backend con Django)}
\label{sec:backend}

El \textit{backend} del \ac{SGDIC} se implementa con \ac{DJ}, un framework
web de alto nivel escrito en Python que provee una arquitectura
de aplicación por componentes. El proyecto se organiza en apps modulares que
agrupan las funcionalidades definidas en los requerimientos del Capítulo~4
del documento de planificación. Cada app contiene su propio paquete de
\texttt{models}, \texttt{views}, \texttt{forms}, \texttt{urls} y pruebas,
y se comunican entre sí a través del \texttt{ORM} de \ac{DJ} y de rutas
\ac{REST} bien definidas.

Las apps principales son:

\begin{description}[leftmargin=2.2cm, labelwidth=2cm]
  \item[usuario] Autenticación, roles (\ac{RBAC}) y recuperación de cuenta.
  \item[cupos] Gestión de disponibilidad, prerrequisitos y prioridad.
  \item[materias] Catálogo, versiones y flujo de aprobación de materias.
  \item[mensajes] Mensajería trazable y notificaciones automáticas.
  \item[evaluaciones] Cuestionarios, evaluacinn docente y registro de faltas.
  \item[quejas] Recepcinn, clasificacinn y trazabilidad de quejas.
  \item[reportes] Tableros y exportaciones para secretarna y departamento.
  \item[analitica] Extraccinn \ac{ETL}, \ac{KPI} y motor de soluciones.
\end{description}

\subsection{Lenguaje y estructura}
\label{sec:backend-lenguaje}

El backend se escribe en \textbf{Python 3.11+} sobre \ac{DJ} 4.x, siguiendo el
patrnn \ac{MVT} (\textit{Model-View-Template}): los \texttt{models} definen el
esquema de la base de datos, las \texttt{views} encapsulan la lngica de
negocio y los \texttt{urls} enrutan las peticiones \ac{REST}. Las vistas
retornan respuesta \ac{JSON} a la API o renderizan plantillas \ac{HTML} al
frontend, segnn corresponda. El \texttt{ORM} de \ac{DJ} traduce
automnticamente las clases del modelo a consultas \ac{SQL}, previniendo
inyeccinn y centralizando las reglas de integridad en migraciones versionadas.

\subsection{Organizacinn por roles}
\label{sec:backend-rol}

El backend filtra la informacinn por rol (\ac{RBAC}) de forma transversal: la
autenticacinn y los permisos se gestionan en la app \texttt{usuario}, y cada
app aplica \texttt{get\_queryset} filtrado por el usuario autenticado. Asn, el
rol \textbf{estudiante} interactna con \texttt{cupos}, \texttt{mensajes} y
\texttt{evaluaciones}; el rol \textbf{docente} aporta a \texttt{materias},
\texttt{evaluaciones} (asistencia, justificaciones, resultados) y
\texttt{mensajes}; y la \textbf{secretarna} accede a \texttt{reportes}
(listas de espera, \ac{KPI}) y \texttt{quejas}, con intervencinn trazada en
hilos escalados. El \textbf{departamento} configura parnmetros globales
(preguntas de evaluacinn, \ac{KPI}) y la app \texttt{analitica} consolida
datos para el motor de soluciones rnpidas.
\begin{figure}[H]
\centering
\adjustbox{max width=\linewidth}{%
\begin{tikzpicture}[
  node distance=1.5cm and 1.5cm,
  nodo/.style={compApp, minimum width=3.4cm, minimum height=1.25cm,
    align=center, inner sep=6pt, font=\small}
]
  \node[nodo, compApi, minimum width=4.0cm] (req) {Petición REST};
  \node[nodo, nodoServidor, below=1.5cm of req, minimum width=5.4cm] (dj) {Servidor de Django};
  \node[nodo, below=1.8cm of dj, xshift=-2.8cm, minimum width=3.3cm] (orm) {ORM de Django};
  \node[nodo, infra, right=2.2cm of orm, minimum width=3.2cm] (ses) {Sesiones / Redis};
  \node[nodo, db, below=1.8cm of orm] (db) {PostgreSQL};
  \node[nodo, infra, below=1.8cm of ses] (notif) {Cola de tareas};
  \draw[flecha] (req) -- (dj);
  \draw[flecha] (dj.south) -- ++(0,-0.75) -| (orm.north);
  \draw[flecha] (dj.south) -- ++(0,-0.75) -| (ses.north);
  \draw[flecha] (orm) -- (db);
  \draw[flecha] (ses) -- (notif);
  \draw[flechaS] (db.east) -- ++(1.1,0) |- (notif.west);
  \node[draw=grisUC, dashed, fit={(orm)(ses)}, inner sep=8pt,
    rounded corners=4pt, label={[font=\small\bfseries]above:WSGI}] {};
\end{tikzpicture}
}%
\caption{Arquitectura del backend de Django del \ac{SGDIC}}
\label{fig:backend}
\end{figure}

\section{Capa de Datos (PostgreSQL)}
\label{sec:datos}

La capa de datos se soporta en PostgreSQL, un sistema gestor de bases de datos
relacionales de código abierto con soporte transaccional \ac{ACID}, tipos de
datos avanzados y alto desempeño en cargas concurrentes. La elección responde a
tres necesidades del \ac{SGDIC}: integridad referencial estricta para el flujo de
cupos y matrículas, consultas analíticas sobre históricos de varios semestres y
capacidad de crecimiento horizontal mediante réplicas de lectura.

El acceso a la base de datos se realiza exclusivamente a través del \ac{ORM} de
\ac{DJ}, que traduce las clases del modelo a tablas y relaciones. Esta decisión
elimina la construcción manual de \ac{SQL}, previene inyección de código y
concentra las reglas de integridad en las migraciones versionadas del proyecto.
Las operaciones de escritura se agrupan en transacciones atómicas
(\texttt{atomic}) para garantizar que un cambio de cupo y su registro de auditoría
se confirmen o se reviertan en conjunto.

La figura~\ref{fig:datos} presenta el modelo físico simplificado. Las entidades
centrales son el usuario y su rol, la asignatura con sus versiones y
prerrequisitos, la oferta de cupos por periodo, la matrícula, la mensajería, la
evaluación y la queja. Las tablas de auditoría y de eventos analíticos cierran el
modelo: la primera registra toda modificación sensible y la segunda alimenta el
repositorio analítico descrito en la Sección~\ref{sec:analitica}.

\begin{figure}[H]
\centering
\adjustbox{max width=\linewidth}{%
\begin{tikzpicture}[
  nodo/.style={db, minimum width=3.2cm, minimum height=1.1cm, align=center,
    inner sep=6pt, font=\small},
  nodol/.style={nodo, minimum width=4.6cm},
  node distance=2.2cm and 2.0cm
]
  \node[nodo] (usuario) {Usuario};
  \node[nodo, right=of usuario] (rol) {Rol};
  \node[nodol, right=of rol] (asignatura) {Asignatura};
  \node[nodo, below=of usuario] (matricula) {Matrícula};
  \node[nodo, below=of rol] (cupo) {Oferta de cupos};
  \node[nodol, below=of asignatura] (version) {Versión y\\[1pt]prerrequisito};
  \node[nodo, below=of matricula] (auditoria) {Auditoría};
  \node[nodo, below=of cupo] (mensaje) {Mensajería};
  \node[nodol, below=of version] (evaluacion) {Evaluación y\\[1pt]queja};
  \draw[flecha] (usuario) -- (rol);
  \draw[flecha] (rol) -- (asignatura);
  \draw[flecha] (usuario) -- (matricula);
  \draw[flecha] (matricula) -- (cupo);
  \draw[flecha] (asignatura) -- (version);
  \draw[flecha] (matricula) -- (auditoria);
  \draw[flecha] (cupo) -- (mensaje);
  \draw[flecha] (version) -- (evaluacion);
\end{tikzpicture}
}%
\caption{Modelo físico simplificado de la base de datos del \ac{SGDIC}}
\label{fig:datos}
\end{figure}

El esquema se administra con migraciones de \ac{DJ} versionadas en el
repositorio, se respalda diariamente con \texttt{pg\_dump} y las consultas más
frecuentes cuentan con índices sobre llaves foráneas, periodos académicos y
estados de trámite.

\section{Comunicación entre Componentes}
\label{sec:comunicacion}

La comunicación sigue el principio de dependencia unidireccional: la capa de
presentación consume a la de aplicación y esta, a la de datos. El navegador
solicita una vista; el servidor de \ac{DJ} la resuelve a través de
\texttt{urls}, \texttt{views} y \texttt{templates}, consulta el \ac{ORM} y
devuelve \ac{HTML} renderizado o un cuerpo \ac{JSON}. Las operaciones que
requieren actualizar la pantalla sin recargar (notificaciones, validación de
prerrequisitos, filtros de cupos) usan llamados asíncronos a la API REST
con \ac{JSON}.

La figura~\ref{fig:secuencia} resume el ciclo completo de una solicitud típica,
desde que el usuario envía el formulario hasta que recibe la respuesta,
incluyendo las tareas que se delegan a la cola de procesos en segundo plano.

\begin{figure}[H]
\centering
\adjustbox{max width=\linewidth}{%
\begin{tikzpicture}[
  node distance=1.9cm and 1.3cm,
  caja/.style={compApp, text width=2.55cm, align=center, minimum height=1.9cm,
    inner sep=6pt, font=\small},
  cajaweb/.style={caja, compWeb},
  cajaapi/.style={caja, compApi},
  cajadb/.style={caja, db},
  etiqueta/.style={font=\tiny, align=center, inner sep=2pt,
    fill=white, fill opacity=0.95, text opacity=1},
  etiquetaAlta/.style={etiqueta, above=7pt},
  etiquetaBaja/.style={etiqueta, below=7pt}
]
  \node[cajaweb] (u) {Frontend\\[2pt]\textit{usuario}};
  \node[cajaapi, right=of u] (v) {Vista\\[2pt]\textit{Django}};
  \node[caja, right=of v] (m) {Modelo\\[2pt]\textit{ORM}};
  \node[cajadb, right=of m] (d) {Datos\\[2pt]\textit{PostgreSQL}};
  \draw[flecha] (u) -- node[etiquetaAlta]{1.~Petición} (v);
  \draw[flecha] (v) -- node[etiquetaAlta]{2.~Consulta} (m);
  \draw[flecha] (m) -- node[etiquetaAlta]{3.~SQL} (d);
  \draw[flechaS] (d) to[out=-50, in=-130, looseness=1.5]
    node[etiquetaBaja, pos=0.5]{4.~Registros} (m);
  \draw[flechaS] (m) to[out=-50, in=-130, looseness=1.5]
    node[etiquetaBaja, pos=0.5]{5.~Objetos} (v);
  \draw[flechaV] (v) to[out=70, in=110, looseness=1.6]
    node[etiquetaAlta, pos=0.5]{6.~Respuesta} (u);
  \node[legenda, below=5.2cm of u, text width=14.0cm, font=\footnotesize,
        align=justify, inner sep=8pt] {
    \textbf{Reglas de la capa de aplicación:} validación de permisos por rol,
    transacción atómica para toda escritura, registro en auditoría y
    notificación asíncrona al interesado mediante la cola de tareas.
  };
\end{tikzpicture}
}%
\caption{Secuencia de comunicación entre las capas del \ac{SGDIC}}
\label{fig:secuencia}
\end{figure}

\section{Capa Analítica y Motor de Soluciones}
\label{sec:analitica}

La capa analítica complementa las tres capas funcionales. Un proceso \ac{ETL}
programado extrae de forma periódica los eventos relevantes (matrículas,
cancelaciones, reclamaciones, respuestas de evaluación), los normaliza y los
carga en un repositorio analítico separado. Sobre ese repositorio se calculan
los \ac{KPI} del departamento y se alimenta el motor de soluciones rápidas, que
sugiere alternativas ante situaciones recurrentes, por ejemplo una asignatura
con demanda superior a la oferta o un incremento sostenido de quejas por
periodo.

El aislamiento del repositorio analítico evita que las consultas de reporte
compitan con las transacciones académicas y permite recalcular indicadores sin
afectar el servicio.

\section{Despliegue y Operación}
\label{sec:despliegue}

El sistema se despliega con el patrón de servidor de aplicaciones \ac{WSGI}
detrás de un servidor web inverso que termina \ac{HTTPS}, sirve los archivos
estáticos y canaliza las peticiones dinámicas. La base de datos PostgreSQL se
ejecuta como servicio independiente con acceso restringido por red. La
figura~\ref{fig:despliegue} presenta la topología propuesta.

\begin{figure}[H]
\centering
\adjustbox{max width=\linewidth}{%
\begin{tikzpicture}[
  node distance=2.3cm and 1.7cm,
  nodo/.style={compApp, text width=2.55cm, align=center, minimum height=2.0cm,
    inner sep=7pt, font=\small},
  nodoweb/.style={nodo, compWeb},
  nodoinfra/.style={nodo, infra},
  nododb/.style={nodo, db},
  nodoapi/.style={nodo, compApi}
]
  \node[nodoweb] (cliente) {Cliente web\\[2pt]\textit{(navegador)}};
  \node[nodoinfra, right=of cliente] (proxy) {Servidor web inverso\\[2pt]\textit{HTTPS y estáticos}};
  \node[nodo, right=of proxy] (app) {Django (WSGI)\\[2pt]\textit{procesos de aplicación}};
  \node[nododb, right=of app] (bd) {PostgreSQL\\[2pt]\textit{datos y respaldo}};
  \node[nodoapi, below=2.6cm of proxy, text width=4.6cm] (cola) {Cola de tareas\\[2pt]\textit{procesos en segundo plano}};
  \draw[flecha] (cliente) -- (proxy);
  \draw[flecha] (proxy) -- (app);
  \draw[flecha] (app) -- (bd);
  \draw[flechaS] (app.south) |- (cola.east);
  \draw[flechaV, dashed] (cola.west) -| (cliente.south);
\end{tikzpicture}
}%
\caption{Topología de despliegue del \ac{SGDIC}}
\label{fig:despliegue}
\end{figure}

La operación prevé integración y despliegue continuos (\ac{CI} y \ac{CD}):\ cada
entrega ejecuta las pruebas automatizadas, aplica las migraciones de base de
datos y publica la versión en el entorno de producción, con posibilidad de
reversión ante fallos. Los respaldos diarios, las migraciones versionadas y los
índices sobre las consultas críticas constituyen el plan de contingencia y
continuidad del servicio.

\section{Seguridad y Control de Acceso}
\label{sec:seguridad}

La seguridad se apoya en cuatro medidas. Primero, autenticación obligatoria con
contraseñas almacenadas mediante funciones de derivación de clave y sesiones
con expiración. Segundo, autorización por rol (\ac{RBAC}) aplicada en el
servidor: estudiante, docente, secretaría y administrador acceden únicamente a
las vistas y operaciones que les corresponden. Tercero, trazabilidad completa:
toda modificación sensible deja registro en la tabla de auditoría, con usuario,
fecha y valor anterior. Cuarto, protección del canal con \ac{HTTPS}, validación
de formularios contra falsificación de petición y control de acceso a los
recursos estáticos.

Estas medidas responden a los requerimientos no funcionales de seguridad,
trazabilidad y disponibilidad del documento de planificación y son condición
previa para la puesta en producción del sistema.





